% Generated draft: requires mathematical and source review.
\documentclass[11pt]{article}
\usepackage[margin=25mm]{geometry}
\usepackage{fontspec}
\setmainfont{Latin Modern Roman}
\usepackage{amsmath,amssymb,mathtools,unicode-math}
\setmathfont{Latin Modern Math}
\usepackage{enumitem,needspace,xurl,hyperref}
\setlength{\parindent}{0pt}
\setlength{\parskip}{5pt}
\setlength{\emergencystretch}{3em}
\newcommand{\N}{\mathbb{N}}
\newcommand{\Z}{\mathbb{Z}}
\newcommand{\Q}{\mathbb{Q}}
\newcommand{\R}{\mathbb{R}}
\newcommand{\C}{\mathbb{C}}
\newcommand{\F}{\mathbb{F}}
\newcommand{\A}{\mathbb{A}}
\newcommand{\PP}{\mathbb{P}}
\newcommand{\E}{\mathbb{E}}
\newcommand{\eps}{\varepsilon}
\newcommand{\dd}{\,\mathrm{d}}
\begin{document}
\section*{UnsolvedMath ID: 30001425}
\textbf{Source problem code:} catalogue entry
\subsection{Definitions and mathematical statement}
For terminology only, consider the standard setting of finite simple undirected graphs. If \(G=(V,E)\) is such a graph and \(A,B\subseteq V\) are disjoint nonempty vertex sets, define
\[
 d_G(A,B)=\frac{|\{\{a,b\}\in E:a\in A,\ b\in B\}|}{|A||B|}.
\]
For \(0<\varepsilon<1\), the pair \((A,B)\) is \(\varepsilon\)-regular if every \(X\subseteq A\) and \(Y\subseteq B\) satisfying \(|X|\geq\varepsilon|A|\) and \(|Y|\geq\varepsilon|B|\) also satisfy
\[
 |d_G(X,Y)-d_G(A,B)|\leq\varepsilon.
\]
These definitions provide standard context; the supplied record does not specify the graph class or the precise version of the regularity lemma intended.

The source denotes by \(M(\varepsilon)\) a bound on the number of parts in regularity partitions. It does not specify the partition convention, an initial lower bound on the number of parts, an exceptional-set convention, or any suppressed parameters. For comparison, define a tower of twos by
\[
 T(1)=2,\qquad T(h+1)=2^{T(h)}\quad(h\in\N,\ h\geq1).
\]
The source describes the partition bounds as tower-type, without specifying their tower height as a function of \(\varepsilon\).

\textbf{Source-level question.} Can applications of the graph regularity lemma be proved, by new proofs, with quantitative estimates substantially better than the estimates obtained through the tower-type partition bound \(M(\varepsilon)\)?

\textbf{Unresolved formulation.} This is a research-program question rather than a specified quantified conjecture. The supplied record does not identify the applications under consideration, their hypotheses and conclusions, the parameters in their quantitative estimates, the benchmark bounds, or the meaning of ``substantially better.'' It also does not determine whether the question concerns some applications, all applications discussed in the report, or another collection. Consequently, no single precise asymptotic inequality or universal quantification over applications can be extracted without additional source information. The question concerns quantitative bounds for applications; it does not explicitly ask for an improved partition bound \(M(\varepsilon)\) itself.
\subsection{Short English statement}
The problem asks whether new proofs of applications of graph regularity can give substantially better quantitative bounds than proofs using tower-type regularity partitions. The applications and the required size of the improvement are unspecified, so the supplied record does not determine a single precise conjecture. Its partially solved label is supplied dated curation.
\subsection{Sources}
\begin{itemize}
\item Local registry: \texttt{data/unsolvedmath/problems.json}, supplied record with permanent UnsolvedMath ID \(30001425\), title ``Quantitative Improvements to Graph Regularity Applications,'' and exact source code \texttt{catalogue entry}. The record attributes the question to Oberwolfach workshop discussions from 2010; no individual proposer is supplied.
\item Supplied report citation: \emph{Graph Theory} (2010), Oberwolfach Reports, pp.~521--580, DOI \texttt{10.4171/owr/2010/11}. \url{https://doi.org/10.4171/owr/2010/11}. This is also the supplied literature-source entry ``Oberwolfach Report 11/2010.'' No exact problem page or report passage is supplied.
\item Supplied literature-source entry: ``Fox, 2011.'' \url{https://doi.org/10.4007/annals.2011.174.1.17}.
\item Supplied literature-source entry: ``Conlon--Fox, 2012.'' \url{https://arxiv.org/abs/1107.4829}.
\item Literature assessment dated exactly \texttt{2026-08-21}: the supplied curation labels the problem partially solved and reports major later quantitative improvements, including removal-lemma bounds and direct proofs, while describing the broad request covering the report's applications as not a single fully settled theorem. This is supplied dated curation, not an independently verified literature assessment; its time-sensitive status and cited claims require independent verification.
\item The supplied statement status is \texttt{corrected\_verified}. Its accompanying verification note says that the mathematical content was checked against the official report and that notation, context, and question form were normalized without changing the source-level claim. That verification is reported as source metadata and has not been independently performed for this draft.
\end{itemize}
\end{document}
